%%%%%%%%%%%%%%%%%%%%%%%%%%%%%%%%%
%
%       EXERCÍCIO
%
%%%%%%%%%%%%%%%%%%%%%%%%%%%%%%%%%

\definevalorquestao

\begin{exercicioBanco}[\valorquestao]
Seja \(p(x) = mx^{2} + nx + 1\). Se \(p(2) = 0\) e \(p(-1) = 0\), então os valores de \(m\) e \(n\) são, respectivamente, iguais a?

% Define as alternativas
\newcommand{\alternativas}{%
    \begin{center}
        \begin{tabularx}{\textwidth}{XXXXX}
            (a) \(-\dfrac{1}{2}\) e \(\dfrac{1}{2}\). &
            (b) \(-1\) e \(1\). &
            (c) \(1\) e \(\dfrac{1}{2}\). &
            (d) \(-1\) e \(-\dfrac{1}{2}\). &
            (e) \(1\) e \(-\dfrac{1}{2}\).
        \end{tabularx}
    \end{center}
}

% Define a resposta correta
\newcommand{\resposta}{A}

% Lógica condicional para exibição
\ifdefstring{\atividade}{lista}{%
    \alternativas
    \vspace{0.5em}
    
    \noindent\textbf{Resposta:} letra \textbf{\resposta}.
}{%
    \ifdefstring{\modo}{objetivo}{%
        \alternativas
    }{%
        % modo = discursiva → não mostra alternativas
    }
}
\end{exercicioBanco}

